\section{DBの開始}

\subsection{実施にかかる手続き（労使合意、承認認可）}


\begin{InyouJobunbox}
  \begin{itemize}
    \item 法第3条
    \item 則第2条〜第4条、第11条
    \item 解釈第8-2、8-7、8-1
  \end{itemize}
\end{InyouJobunbox}


\begin{fillblank}[自作 2026/10/4]

  \noindent \textbf{◯DB法}
  \begin{lawstructure}
    \begin{lawarticle}{第三条}{確定給付企業年金の実施}
      厚生年金適用事業所の事業主は、確定給付企業年金を実施しようとするときは、確定給付企業年金を実施しようとする厚生年金適用事業所に使用される（　　A　　）で組織する労働組合があるときは当該労働組合、当該（　　A　　）で組織する労働組合がないときは当該（　　A　　）を代表する者の同意を得て、確定給付企業年金に係る規約（以下「規約」という。）を作成し、次の各号のいずれかに掲げる手続を執らなければならない。
      \lawitem{一}当該規約について厚生労働大臣の承認を受けること。
      \lawitem{二}企業年金基金（以下「基金」という。）の設立について厚生労働大臣の認可を受けること。

      \lawparagraph{2}
      確定給付企業年金は、一の厚生年金適用事業所について一に限り実施することができる。ただし、政令で定める場合においては、この限りでない。
      \lawparagraph{3}
      二以上の厚生年金適用事業所について確定給付企業年金を実施しようとする場合においては、第一項の同意は、（　　B　　）について得なければならない。
    \end{lawarticle}
  \end{lawstructure}

\end{fillblank}

\begin{answerbox}
 \begin{enumerate}[A:]
    \item 厚生年金保険の被保険者の過半数
    \item 各厚生年金適用事業所
  \end{enumerate}
\end{answerbox}


\begin{fillblank}[自作 2026/10/4]

  \noindent \textbf{◯DB則}
  \begin{lawstructure}
    \begin{lawarticle}{第三条}{過半数代表者}
      法第三条第一項、法第六条第二項（法第七条第二項において準用する場合を含む。）及び法第七十八条第一項の規定による手続を厚生年金保険の被保険者（法第二条第三項に規定する厚生年金保険の被保険者をいう。以下同じ。）の過半数を代表する者（以下この条において「過半数代表者」という。）の同意を得て行う場合にあっては、当該過半数代表者は、次の各号のいずれにも該当する者とする。
      \lawitem{一}労働基準法（昭和二十二年法律第四十九号）第四十一条第二号に規定する（　　A　　）の地位にある者でないこと。
      \lawitem{二}過半数代表者を選出することを明らかにして実施される（　　B　　）等の方法による手続により選出された者であって、（　　C　　）に基づき選出されたものでないこと。
      
      \lawparagraph{2}
      前項第（　　※出題の都合で略　　）号に該当する者がいない厚生年金適用事業所にあっては、過半数代表者は同項（　　D　　）に該当する者とする。
      \lawparagraph{3}
      確定給付企業年金を実施しようとする又は実施する厚生年金適用事業所の事業主は、当該事業主に使用される者が過半数代表者であること若しくは過半数代表者になろうとしたこと又は過半数代表者として正当な行為をしたことを理由として（　　E　　）をしないようにしなければならない。
      \lawparagraph{4〜5} （略）
    \end{lawarticle}
  \end{lawstructure}

\end{fillblank}

\begin{answerbox}
 \begin{enumerate}[A:]
    \item 監督又は管理
    \item 投票、挙手
    \item 事業主の意向
    \item 第二号
    \item 不利益な取り扱い
  \end{enumerate}
\end{answerbox}



\begin{truefalse}[年金1 \ 2019 \ 問題1(2)\textcircled{6}]
\textbf{下線部の内容の正誤を判定し、誤っている場合は正しい内容に改めなさい。}
\\

確定給付企業年金法施行規則第3 条において、規約の作成、規約の変更および実施事業所
の増減等に関する同意手続きを、厚生年金保険の被保険者の過半数を代表する者の同意を得
て行う場合には、当該者は次の A、B のいずれにも該当する者とすることとされている。
\begin{enumerate}[A]
  \item 労働基準法第41 条第2 号に規定する監督または管理の地位にある者でないこと
  \item 過半数代表者を選出することを明らかにして実施される投票、挙手等の方法による手続
きにより選出された者\underline{であること}
\end{enumerate}
\end{truefalse}

\begin{answerbox}
 \textbf{誤} \\
  〜により選出された者\underline{であって、事業主の意向により選出された者でないこと}
\end{answerbox}

\begin{truefalse}[自作 \ 2026/10/4]
\textbf{以下の文章の正誤を判定し、誤っている場合はその理由を記述しなさい。}
\\

厚生年金適用事業所の事業主は、確定給付企業年金を実施しようとするとき、
確定給付企業年金を実施しようとする厚生年金適用事業所に使用される厚生年金保険の被保険者の過半数で
組織する労働組合があるときであっても、
当該厚生年金保険の被保険者の過半数を代表する者の同意さえ得られれば、確定給付企業年金を実施できる。
\end{truefalse}

\begin{answerbox}
 \textbf{誤} \\
  「過半数組合の同意」か「過半数代表者の同意」かの選択は、過半数組合の有無によって自動で決まるものであり、
  事業主が任意に選択できるものではない。
\end{answerbox}

\begin{explanationbox}
  出典：DB法第三条
\end{explanationbox}

\newpage

\subsection{１事業所１DBの原則}



\begin{InyouJobunbox}
  \begin{itemize}
    \item 法第3条第2項
    \item 令第1条
    \item 則第1条
  \end{itemize}
\end{InyouJobunbox}


\begin{fillblank}[自作 2026/10/4]

  \noindent \textbf{◯DB令}
  \begin{lawstructure}
    \begin{lawarticle}{第一条}{複数の確定給付企業年金を実施できる場合}
      確定給付企業年金法（以下「法」という。）第三条第二項ただし書の政令で定める場合は、一の厚生年金適用事業所（法第二条第二項に規定する厚生年金適用事業所をいう。以下同じ。）について二の確定給付企業年金を実施する場合であって当該二の確定給付企業年金のうちいずれか一方の確定給付企業年金を実施する厚生年金適用事業所（以下「実施事業所」という。）の事業主（第五十三条並びに附則第三条及び第八条を除き、以下「事業主」という。）の全部が同時に（　　A　　）とならないときその他厚生労働省令で定める場合とする。
    \end{lawarticle}
  \end{lawstructure}

  \noindent \textbf{◯DB則}
  \begin{lawstructure}
    \begin{lawarticle}{第一条}{複数の確定給付企業年金を実施できる場合}

      確定給付企業年金法施行令（平成十三年政令第四百二十四号。以下「令」という。）第一条の厚生労働省令で定める場合は、次のとおりとする。

      \lawitem{一}
      一の厚生年金適用事業所（確定給付企業年金法（平成十三年法律第五十号。以下「法」という。）第二条第二項に規定する厚生年金適用事業所をいう。以下同じ。）について二以上の確定給付企業年金を実施する場合であって、それぞれの確定給付企業年金の加入者（以下「加入者」という。）について適用される（　　B　　）、（　　C　　）その他これらに準ずるもの（以下「（　　B　　）等」という。）が異なる場合
      \lawitem{二}
      法人である確定給付企業年金を実施する事業主（第三条第一項第二号、第三項及び第五項、第十九条の二第二号イ、第百二十条、附則第六条第一項第一号、附則第七条第一項並びに附則第十二条第一項第一号を除き、以下「事業主」という。）が他の法人である事業主と（　　D　　）した場合であって、当該（　　D　　）の日から起算して（　　E　　）を経過していない場合
      \lawitem{三}
      給付の額の算定方法が第二十五条第四号に掲げる方法である確定給付企業年金（以下「（　　F　　）」という。）と（　　F　　）でない確定給付企業年金とをそれぞれ実施する場合
    \end{lawarticle}
  \end{lawstructure}

\end{fillblank}


\begin{answerbox}
 \begin{enumerate}[A:]
    \item 他方の確定給付企業年金の事業主の全部
    \item 労働協約
    \item 就業規則
    \item 合併
    \item 原則として一年
    \item リスク分担型企業年金
  \end{enumerate}
\end{answerbox}


\begin{truefalse}[年金1 \ 2019 \ 問題1(2)\textcircled{3}]
\textbf{下線部の内容の正誤を判定し、誤っている場合は正しい内容に改めなさい。}
\\

例外的に2 つ以上の確定給付企業年金を実施することが認められているケースとして、会
社合併した場合であって、\underline{当該合併にかかる規約変更の日から起算して2 年}を経過していな
い場合が該当する。
\end{truefalse}

\begin{answerbox}
 \textbf{誤} \\
  〜合併した場合であって、\underline{当該合併の日から起算して原則として1年}を経過していない場合〜
\end{answerbox}



\begin{truefalse}[年金1 \ H29 \ 問題2(1)\textcircled{3}]
\textbf{以下の文章の正誤を判定し、誤っている場合その理由を記述しなさい。}
\\

リスク分担型企業年金とリスク分担型企業年金でない確定給付企業年金の経理をそれぞれで
行うことを規約に定めることで、１つの規約型確定給付企業年金において、リスク分担型企業
年金とリスク分担型企業年金でない確定給付企業年金を併存することができる。
\end{truefalse}

\begin{answerbox}
 \textbf{誤} \\
  経理をそれぞれ行うとともに、資産をそれぞれ区分して運用することを規約に定める必要がある。
\end{answerbox}


\begin{explanationbox}
  参照：承認認可基準別紙第3-2(4) 

  \textgt{
  リスク分担型企業年金とリスク分担型企業年金でない確定給付企業年金の経理をそれぞれで行うとともに、資産をそれぞれに区分して運用することを規約に定める措置。なお、基金型の場合においては、2-4(6)の事項を規約に定める措置もあわせて講じること。
  }
\end{explanationbox}



\begin{shortanswer}[年金1 \ 2023 \ 問題2(1)ア]
１つの厚生年金適用事業所が複数の確定給付企業年金を実施することができる場合として、
確定給付企業年金法施行令第１条に定める場合（注）\underline{以外}の、確定給付企業年金法施行規則
第１条に定める場合を３つ簡潔に入力しなさい。 （250 字以内）\\
（注）複数の厚生年金適用事業所が共同で実施する確定給付企業年金に加入する一方で、企業
独自に実施する確定給付企業年金に加入する等、企業年金を実施する事業主の範囲が異
なる場合。
\end{shortanswer}

\begin{answerbox}
 \begin{itemize}
  \item それぞれの確定給付企業年金の加入者について適用する労働協約、就業規則その他これらに準ずるものが異なる場合
  \item 法人である確定給付企業年金を実施する事業主が他の法人である事業主と合併した場合であって、当該合併した日から起算して原則として一年を経過していない場合
  \item リスク分担型企業年金とリスク分担型企業年金でない確定給付企業年金とをそれぞれ実施する場合
 \end{itemize}
\end{answerbox}

\begin{explanationbox}
  類題：年金2 2022 問題2(3)
\end{explanationbox}

\newpage

\subsection{規約に定める事項}



\begin{InyouJobunbox}
  \begin{itemize}
    \item 法第4条、第11条
    \item 令第2条、第5条
    \item 則第14条
    \item 解釈第2
  \end{itemize}
\end{InyouJobunbox}


\begin{fillblank}[年金2 \ 2018 \ 問題1(3)A]

  \noindent \textbf{◯DB令}
  \begin{lawstructure}
    \begin{lawarticle}{第二条}{規約型企業年金の規約で定めるその他の事項}
      法第四条第九号の政令で定める事項は、次のとおりとする。

      （略）

      \lawitem{四}
      法第八十一条の二第二項、第八十二条の六第一項又は第九十一条の二十七第二項の規定に基づき、当該確定給付企業年金の資産管理運用機関等（法第三十条第三項に規定する資産管理運用機関等をいう。以下同じ。）が脱退一時金相当額（法第八十一条の二第一項に規定する脱退一時金相当額をいう。以下同じ。）若しくは積立金（法第五十九条に規定する積立金をいう。以下同じ。）、個人別管理資産額（確定拠出年金法（平成十三年法律第八十八号）第二条第十三項に規定する個人別管理資産額をいう。以下この号において同じ。）又は中小企業退職金共済法（昭和三十四年法律第百六十号）第十七条第一項に規定する厚生労働省令で定める金額若しくは同法第三十一条の四第一項に規定する（　　A　　）に相当する額の移換又は引渡しを受ける場合にあっては、当該脱退一時金相当額若しくは積立金、個人別管理資産額又は同法第十七条第一項に規定する厚生労働省令で定める金額若しくは同法第三十一条の四第一項に規定する（　　A　　）に相当する額の移換又は引渡しに関する事項
      
      （略）
    \end{lawarticle}
  \end{lawstructure}

\end{fillblank}

\begin{answerbox}
 \begin{enumerate}[A:]
  \item 解約手当金
 \end{enumerate}
\end{answerbox}

\newpage

\subsection{規約の承認認可基準}



\begin{InyouJobunbox}
  \begin{itemize}
    \item 法第5条、第12条〜第15条
    \item 令第3条〜第4条、第6条〜第10条
  \end{itemize}
\end{InyouJobunbox}

\begin{fillblank}[自作 2026/10/7]

  \noindent \textbf{◯DB法}
  \begin{lawstructure}
    \begin{lawarticle}{第十二条}{基金の設立認可の基準等}
      厚生労働大臣は、第三条第一項第二号の設立の認可の申請があった場合において、当該申請が次に掲げる要件に適合すると認めるときは、同号の認可をするものとする。

      （略）

      \lawitem{四}
      当該申請に係る事業所において、常時政令で定める数以上の加入者となるべき厚生年金保険の被保険者を使用していること、又は使用すると見込まれること（次号に掲げる場合を除く。）。

      \lawitem{五}
      厚生年金適用事業所の事業主が共同して基金を設立しようとする場合にあっては、当該事業主の当該申請に係る事業所において、（　　A　　）して、常時政令で定める数以上の加入者となるべき厚生年金保険の被保険者を使用していること、又は使用すると見込まれること。

      （略）
    \end{lawarticle}

  \noindent \textbf{◯DB令}

    \begin{lawarticle}{第六条}{基金の設立に必要な厚生年金保険の被保険者の数}
      法第十二条第一項第四号及び第五号の政令で定める数は、（　　B　　）人とする。
    \end{lawarticle}

  \end{lawstructure}

\end{fillblank}

\begin{answerbox}
 \begin{enumerate}[A:]
  \item 合算
  \item 300
 \end{enumerate}
\end{answerbox}

\newpage


\subsection{規約型DBの規約変更（給付減額以外）}


\begin{InyouJobunbox}
  \begin{itemize}
    \item 法第6条〜第7条
    \item 則第7条〜第10条
    \item 解釈第1-3
  \end{itemize}
\end{InyouJobunbox}

\begin{truefalse}[年金1 \ 2019 \ 問題1(2)\textcircled{4} \ (改題)]
次のA～E の5 つの規約型企業年金の規約変更のうち、確定給付企業年金法施行規則第10
条に定める、厚生労働大臣の認可を受ける必要のない軽微な規約変更で届出の必要のないも
のを全て挙げなさい。

\begin{enumerate}[A:]
  \item 実施事業所の事業主の住所の変更（市町村の名称の変更に伴う場合）
  \item リスク分担型企業年金の調整率
  \item 毎事業年度の特別掛金の額に関する事項のうち弾力償却または定率償却による毎事業年度の掛金の額に関する事項の変更
  \item 確定給付企業年金の移転承継に関する企業年金の名称の変更
  \item 法令の改正に伴う変更に係る事項の変更（給付の実質的な変更を伴うものを除く）
\end{enumerate}

\end{truefalse}

\begin{answerbox}
 A, B, C, E
\end{answerbox}


\begin{explanationbox}
  \begin{enumerate}[A:]
    \item 届出不要 \& 労使合意不要
    \item 届出不要 \& 労使合意不要
    \item 届出不要 \& 労使合意要
    \item 届出 \qquad  \& 労使合意要
    \item 届出不要 \& 労使合意不要
  \end{enumerate}
\end{explanationbox}


\begin{shortanswer}[年金1 \ 2021 \ 問題2(1)\textcircled{1}]
確定給付企業年金法施行規則および『確定給付企業年金法並びにこれに基づく政令及び省令
について（法令解釈）』の改正により、2020 年（令和2 年）10 月1 日付で、確定給付企業年金
法第6 条第1 項の「厚生労働省令で定める軽微な変更」に追加された規約変更内容を２つ簡記
しなさい。
\end{shortanswer}

\begin{explanationbox}
  昔の時事問題。

  ヒント\textcircled{1}：繰り下げ規定の変更

  ヒント\textcircled{2}：リスク対応掛金についての変更
\end{explanationbox}


\begin{answerbox}
 \begin{itemize}
  \item 老齢給付金または脱退一時金の支給の繰り下げ規定の新設
  \item 少なくとも5年ごとに行う財政再計算において、対応後リスク充足額が財政悪化リスク相当額を
    上回ることとなる場合に行うリスク対応掛金額の減少又は拠出の終了
 \end{itemize}
\end{answerbox}

\begin{comment}
  

\begin{shortanswer}[自作 \ 2026/10/7]
確定給付企業年金法施行規則および『確定給付企業年金法並びにこれに基づく政令及び省令
について（法令解釈）』の改正により2024年12月1日付で、確定給付企業年金法第6条第1項の
「承認」を受けるべき規約変更が追加された。その規約変更内容を簡記しなさい。
\end{shortanswer}

\begin{answerbox}
加入者の過去分の給付増額を行う場合であって、当該増額にかかる実施事業所の事業主が企業型DCも実施している場合
\end{answerbox}

\end{comment}

\begin{shortanswer}[年金1 \ 2024 \ 問題2(1)ウ]
確定給付企業年金の規約の変更において厚生労働大臣の承認・認可を受ける必要のない軽微な
変更として、確定給付企業年金法第４条第５号に掲げる事項に係る変更のうち軽微な変更があ
る。給付の額を減額する場合は当該軽微な変更に該当しないが、2024 年（令和 6 年）12 月 1
日以降、当該軽微な変更に該当しないとされる場合が追加された。当該追加された内容につい
て簡潔に入力しなさい。なお、解答にあたっては通知『確定給付企業年金法並びにこれに基づ
く政令及び省令について（法令解釈）』 第１の３の内容も踏まえて記載すること。 （250 字以内）
\end{shortanswer}


\begin{answerbox}
規約の変更が効力を有することとなる日前の期間に係る給付の額の増額（当該増額に係る実施
事業所の事業主が企業型年金を実施している場合に限る。）に該当する場合。

なお、当該増額に係る実施事業所の事業主が企業型年金を実施している場合とは、当該給付の
額が増額されることとなる加入者等が企業型年金加入者である場合をいう。
\end{answerbox}

\begin{explanationbox}
  則第7条第1項第4号の以下の部分。

  \begin{lawstructure}
    \begin{lawarticle}{第七条}{規約の軽微な変更等}
      法第六条第一項の厚生労働省令で定める軽微な変更は、次に掲げる事項の変更とする。
      
      （略）
      \lawitem{四}
      法第四条第五号に掲げる事項（労働協約等の変更により法第二十七条の規定による
      加入者の資格の喪失の時期が変更になる場合その他の給付の設計の軽微な変更
      （\textcolor{red}{給付の額を減額する場合及び規約の変更が効力を有することとなる日（第八十五条の三において「規約変更日」という。）前の期間に係る給付の額を増額する場合（当該増額に係る実施事業所の事業主が企業型年金を実施している場合に限る。）を除く。}）がある場合に限り、第九号に掲げる事項を除く。）
      
      （略）
    \end{lawarticle}
  \end{lawstructure}
\end{explanationbox}

\newpage

\subsection{基金型DBの規約変更（給付減額以外）}


\begin{InyouJobunbox}
  \begin{itemize}
    \item 法第16条〜第17条
    \item 則第15条〜第18条
  \end{itemize}
\end{InyouJobunbox}

出題例なし

\newpage

\subsection{給付減額にかかる規約変更（リスク分担型DB以外）}

\begin{InyouJobunbox}
  \begin{itemize}
    \item 令第4条、第7条
    \item 則第5条、第6条、第12条、第13条
    \item 解釈1-2(1)\textcircled{1}\textcircled{2}\textcircled{4}\textcircled{5}\textcircled{7}〜\textcircled{10}、解釈1-2(2)、解釈1-2(3)\textcircled{1}
  \end{itemize}
\end{InyouJobunbox}


\subsubsection{給付減額の理由}

\begin{fillblank}[年金1 \ 2022 \ 問題1(1)アイ]

  \noindent \textbf{◯法令解釈}
  \begin{lawstructure}
    \begin{lawarticle}{}{}
      \textbf{第１ 規約の承認又は基金の設立認可の基準に関する事項}
      \lawparagraph{1} （略）
      \lawparagraph{2} 
      給付の額を減額する場合の取扱い
      \lawitem{(1)}
      給付の額を減額する場合にあっては、次に掲げる事項について留意すること。
        \hangedlawitem{\textcircled{1}〜\textcircled{4}　（略）}{3}
        
        \hangedlawitem{\textcircled{5}}{3}
        給付設計の変更日における受給権者等（加入者である受給権者及び加入者で
あった者をいう。以下同じ。）の給付の額は、原則として引下げの対象とすべ
きではなく、仮に引き下げる場合でも、（　　A　　）するために真
にやむを得ない場合に限り行われるものであること。この場合においては次の
措置を講じる必要があること。
        \hangedlawitem{ア}{4}
        事業主、加入者及び受給権者等の三者による協議の場を設けるなど受給権
者等の（　　B　　）させる措置を講じること。
        \hangedlawitem{イ}{4}
        全受給権者等に対し、事前に、給付設計の変更に関する十分な説明と意向
確認を行っていること。
        \hangedlawitem{\textcircled{6}〜\textcircled{10}　（略）}{3}

    \end{lawarticle}
  \end{lawstructure}

\end{fillblank}


\begin{answerbox}
  \begin{enumerate}[A:]
    \item 確定給付企業年金を存続
    \item 意向を十分に反映
  \end{enumerate}
\end{answerbox}

\begin{truefalse}[年金1 \ 2023 \ 問題1(3)(ウ)]
  \textbf{下線部の正誤を判定し、誤っている場合は正しい内容に改めなさい。}
\\

  「給付の額を減額しない場合に増加する掛金の額が\underline{事業主の当期純利益の過去３年間程度の平均の概ね２割以上}
  となっている」場合、確定給付企業年金法施行規則第５条第２号に定めら
れている給付減額の理由に該当する。
\end{truefalse}

\begin{answerbox}
  \textbf{誤}

  事業主の当期純利益の過去5年間程度の平均の概ね1割以上
\end{answerbox}

\begin{truefalse}[年金1 \ 2018 \ 問題2(1)\textcircled{3}]
  \textbf{以下の内容の正誤を判定し、誤っている場合はその理由を記述しなさい。}
  \\

  過去3年間程度の平均において、実施事業所の事業主の当期純利益がマイナスまたはその
見込みである場合には、確定給付企業年金法施行規則第5条第2号の「事業主が掛金を拠出
することが困難になると見込まれるため、給付の額を減額することがやむを得ないこと」に
該当する。
\end{truefalse}

\begin{answerbox}
  \textbf{誤}

  「過去3年間程度の平均」ではなく、「過去5年間程度の過半数の期」
\end{answerbox}

\begin{fillblank}[数理人試験 \ 2019 \ 問題4 \ 設問1]
  \noindent \textbf{◯DB則}
  \begin{lawstructure}
    \begin{lawarticle}{第五条}{給付減額の理由}
        令第四条第二号の厚生労働省令で定める理由は、次のとおりとする。
        ただし、加入者である受給権者（給付を受ける権利（以下「受給権」という。）を有する者をいう。
        以下同じ。）及び加入者であった者（以下「受給権者等」という。）の給付
        （加入者である受給権者にあっては、当該受給権に係る給付に限る。）
        の額を減額する場合にあっては、（　　A　　）に掲げる理由とする。

        \lawitem{一}
        確定給付企業年金を実施する厚生年金適用事業所（以下「実施事業所」という。）において労働協約等が変更され、
        その変更に基づき給付の設計の見直し（略）を行う必要があること。\\
        （以下略）

        \lawitem{二}
        実施事業所の経営状況の悪化又は（　　B　　）により、
        事業主が掛金を拠出することが困難になると見込まれるため、
        給付の額を減額すること（リスク分担型企業年金開始変更、
        リスク分担型企業年金終了変更又はリスク分担型企業年金統合等変更を行った結果、
        給付の額が減額されることとなる場合を含む。次号において同じ。）がやむを得ないこと。

        \lawitem{三}
        法第七十四条第一項の規定により規約型企業年金を他の規約型企業年金と統合する場合、法第七十九条第二項又は第八十一条第二項の規定により事業主が給付の支給に関する権利義務を承継する場合であって、給付の額を減額することにつきやむを得ない事由があること。

        \lawitem{四}
        給付の額を減額し、当該事業主が拠出する掛金のうち給付の額の減額に伴い減少する額に
        相当する額を（　　C　　）（確定拠出年金法（平成十三年法律第八十八号）第三条第三項第七号に
        規定する（　　C　　）をいう。）に充てること又は法第八十二条の二第一項の規定により、
        給付に充てるべき積立金（以下「積立金」という。）の一部を、
        実施事業所の事業主が実施する企業型年金（確定拠出年金法第二条第二項に規定する企業型年金
        をいう。以下同じ。）の（　　D　　）（同条第七項第一号ロに規定する（　　D　　）をいう。
        以下同じ。）に移換すること。

        \lawitem{五}
        当該規約の変更がリスク分担型企業年金開始変更を内容とするものである場合において、変更後のリスク分担型企業年金が第二十五条の二第一項第二号イに規定する場合に該当することとなること又は該当することとなる蓋然性が高いこと。

        \lawitem{六}
        当該規約の変更がリスク分担型企業年金終了変更を内容とするものである場合において、変更前のリスク分担型企業年金が第二十五条の二第一項第二号ロに規定する場合に該当していること又は該当する蓋然性が高いこと。
    \end{lawarticle}
  \end{lawstructure}
\end{fillblank}

\begin{answerbox}
  \begin{enumerate}[A:]
    \item 第二号、第五号及び第六号
    \item 掛金の額の大幅な上昇
    \item 事業主掛金
    \item 資産管理機関
  \end{enumerate}
\end{answerbox}

\begin{shortanswer}[数理人試験 \ 2019 \ 問題4 \ 設問2]
  確定給付企業年金法施行規則第五条第三号の「やむを得ないこと」 の内容について、通知「確
定給付企業年金制度について」に記載されている内容を簡記せよ。
\end{shortanswer}

\begin{answerbox}
  合併等により給付設計の変更を行わなければ給付水準に大幅な格差が生じることとなるため、当該格差を是正する必要があること。
\end{answerbox}

\begin{explanationbox}
  類題：年金1 2021 問題2(1)\textcircled{3}
\end{explanationbox}

\begin{shortanswer}[数理人試験 \ 2019 \ 問題4 \ 設問3]
  確定給付企業年金法施行規則第五条第五号の「該当することとなる蓋然性が高いこと」の内
容について、通知「確定給付企業年金制度について」に記載されている内容を簡記せよ。
\end{shortanswer}

\begin{answerbox}
  変更後のリスク分担型企業年金において、リスク充足額が財政悪化リスク相当額の2分の1以上であること。
\end{answerbox}



\begin{fillblank}[数理人試験 \ H28 \ 問題1 \ 設問1（改題）]
  \noindent \textbf{◯法令解釈}
  \begin{lawstructure}
    \begin{lawarticle}{}{}
      \lawparagraph{2} 
      給付の額を減額する場合の取扱い
      \lawitem{(1)}
      給付の額を減額する場合にあっては、次に掲げる事項について留意すること。
        \hangedlawitem{\textcircled{1}}{3}
        確定給付企業年金法施行規則（平成１４年厚生労働省令第２２号。以下「規
        則」という。）第５条第２号の「掛金の額の大幅な上昇により、事業主が掛金
        を拠出することが困難と見込まれるため、給付の額を減額することがやむを得
        ない」ことにより給付の額を減額する場合において、確定給付企業年金につい
        て（　　A　　）の規約変更を行っている場合には、当該規約変更時から原則として
        （　　B　　）が経過していること。なお、次のアからウのいずれかに該当する場合には
        同号に該当するものとして取り扱うこと。

        \hangedlawitem{ア}{4}
        過去（　　B　　）間程度のうち過半数の期において、実施事業所の事業主（以下こ
        の\textcircled{1}において「事業主」という。）の（　　C　　）がマイナス又はその見込み
        であること。

        \hangedlawitem{イ}{4}
        給付の額を減額しない場合に増加する掛金の額が事業主の（　　C　　）の過
        去（　　B　　）間程度の平均の概ね（　　D　　）となっていること。
        
        \hangedlawitem{ウ}{4}
        複数の事業主で確定給付企業年金を実施している場合については、アに該
        当する事業主が全事業主の概ね（　　E　　）又はイに該当する事業主が全事業主
        の概ね（　　F　　）となっていること。ただし、一部の事業主が連結決算を行っ
        ている場合は、当該事業主を一の事業主として、当該事業主の増加する掛金
        の額の合計及び連結決算における（　　C　　）を用いることができること。


        \hangedlawitem{\textcircled{2}}{3}
        規則第５条第３号の「やむを得ない事由があること」とは、（　　G　　）により給
        付設計の変更を行わなければ給付水準に大幅な格差が生じることとなるため、
        当該格差を是正する必要がある場合をいうこと（規則第１２条第３号及び規則
        附則第５条第１項の「やむを得ない」も同様。）。

        \hangedlawitem{\textcircled{3}}{3}
        （略）

        \hangedlawitem{\textcircled{4}}{3}
        給付設計の変更日における加入者に対して、（　　H　　）するための適切な
        経過措置を講じること。経過措置を講じることが困難な場合にあっては、その
        旨を加入者に十分に説明した上で、給付の額を減額するものであること。


        （以下略）

    \end{lawarticle}
  \end{lawstructure}
\end{fillblank}

\begin{answerbox}
  \begin{enumerate}[A:]
    \item 給付改善
    \item 5年
    \item 当期純利益
    \item 1割以上
    \item 5割以上
    \item 2割以上
    \item 合併等
    \item 受給権を保全
  \end{enumerate}
\end{answerbox}

\begin{fillblank}[数理人試験 \ 2024 \ 問題1 \ 設問2]
  \begin{lawstructure}
    \noindent \textbf{◯DB則}
    \begin{lawarticle}{第五条}{給付減額の理由}
        令第四条第二号の厚生労働省令で定める理由は、次のとおりとする。
        ただし、加入者である受給権者（給付を受ける権利（以下「受給権」という。）を有する者をいう。
        以下同じ。）及び加入者であった者（以下「受給権者等」という。）の給付
        （加入者である受給権者にあっては、当該受給権に係る給付に限る。）
        の額を減額する場合にあっては、第（　　A　　）号、第五号及び第六号に掲げる理由とする。

        \lawitem{一}
        確定給付企業年金を実施する厚生年金適用事業所（以下「実施事業所」という。）において
        （　　B　　）が変更され、その変更に基づき（　　C　　）の見直し（リスク分担型企業年金
        でない確定給付企業年金をリスク分担型企業年金に変更すること
        （次号及び第五号並びに第十二条第一号及び第二号において「リスク分担型企業年金開始変更」
        という。）、リスク分担型企業年金をリスク分担型企業年金でない確定給付企業年金に変更すること
        （次号及び第六号並びに第十二条第一号及び第二号において「リスク分担型企業年金終了変更」
        という。）及び次に掲げる事由によりリスク分担型企業年金に係る見直しを行うこと
        （次号において「リスク分担型企業年金統合等変更」という。）を含む。）を行う必要があること。
        
        （中略）

        \lawitem{二}
        実施事業所の（　　D　　）の悪化又は掛金の額の（　　E　　）により、
        事業主が掛金を拠出することが困難になると見込まれるため、
        給付の額を減額すること（リスク分担型企業年金開始変更、
        リスク分担型企業年金終了変更又はリスク分担型企業年金統合等変更を行った結果、
        給付の額が減額されることとなる場合を含む。次号において同じ。）がやむを得ないこと。

        \lawitem{三}
        法第七十四条第一項の規定により規約型企業年金を他の規約型企業年金と統合する場合、法第七十九条第二項又は第八十一条第二項の規定により事業主が給付の支給に関する権利義務を承継する場合であって、給付の額を減額することにつきやむを得ない事由があること。

        \lawitem{四}
        給付の額を減額し、当該事業主が拠出する掛金のうち給付の額の減額に伴い減少する額に
        相当する額を（　　F　　）（確定拠出年金法（平成十三年法律第八十八号）第三条第三項第七号に
        規定する（　　F　　）をいう。）に充てること又は法第八十二条の二第一項の規定により、
        給付に充てるべき積立金（以下「積立金」という。）の一部を、
        実施事業所の事業主が実施する企業型年金（確定拠出年金法第二条第二項に規定する企業型年金
        をいう。以下同じ。）の（　　G　　）（同条第七項第一号ロに規定する（　　G　　）をいう。
        以下同じ。）に移換すること。

        （以下略）
    \end{lawarticle}
  \end{lawstructure}
\end{fillblank}

\begin{answerbox}
  \begin{enumerate}[A:]
    \item 二
    \item 労働協約等
    \item 給付設計
    \item 経営状況
    \item 大幅な上昇
    \item 事業主掛金
    \item 資産管理機関
  \end{enumerate}  
\end{answerbox}

\subsubsection{給付減額にかかる規約変更の手続き}

\begin{shortanswer}[年金1 \  2022 \ 問題2(1)イ]
  確定給付企業年金法施行規則第6 条第1 項第2 号において、受給権者等の給付の額を減額す
る場合の手続きについて定められているが、受給権者等のうち希望する者に対する「その他
の当該最低積立基準額が確保される措置」として、『確定給付企業年金法並びにこれに基づく
政令及び省令について（法令解釈）』に例示されているものを3 つ簡潔に入力しなさい。
\end{shortanswer}

\begin{answerbox}
  \begin{enumerate}[(1)]
    \item 規約の変更による給付の額の減額がないものとして算定した最低積立基準額を一時金として支
給する措置に加えて、次の a 又は b その他の給付の額の減額がないものとして合理的に算定し
た額を一時金として支給する選択肢を追加する方法。
    \begin{enumerate}[a]
      \item 規約の変更による給付の額の減額がないものとして、確定給付企業年金法施行規則第 26 条第
3 項に規定する予定利率及び予定死亡率により算定される給付に要する費用の予想額の現価
に相当する額（以下「通常予測給付現価相当額」という。）
      \item 規約の変更による給付の額の減額がないものとして、規約の定めるところにより算定される
一時金として支給する老齢給付金の額（以下「選択一時金の額」という。
    \end{enumerate}
    \item 規約の変更による給付の額の減額がないものとして算定した最低積立基準額から当該規約の変
更による給付に相当する最低積立基準額を控除した額を一時金として支給し、かつ、当該規約
の変更による給付を支給する方法。
    \item 前記(2)の措置に加えて、前記(2)中「最低積立基準額」を「通常予測給付現価相当額又は選択一時
金の額その他合理的に算定した一時金の額」と読み替えて適用する選択肢を追加する方法。
  \end{enumerate}
\end{answerbox}

\begin{shortanswer}[合スト \ DB編上(2026年版) \ p.163]
  給付減額の同意要件を減額対象が加入者、受給権者等の場合でそれぞれ簡記せよ。
\end{shortanswer}

\begin{answerbox}
  ＜加入者＞

  \begin{itemize}
    \item 減額対象者の3分の1以上で組織する労働組合がある場合は、当該労働組合の同意
    \item 減額対象者の3分の2以上の個別同意（減額対象者の3分の2以上で組織する労働組合がある場合は、当該労働組合の同意で代替可能）
    \item 実施事業所が複数ある場合は、実施事業所ごとに上記の同意が必要
  \end{itemize}
  
  ＜受給権者＞

  \begin{itemize}
    \item 減額対象者の3分の2以上の同意
  \end{itemize}
\end{answerbox}

\begin{shortanswer}[年金1 \ H19 \ 問題2(1)（改題）]
  確定給付企業年金の一部を確定拠出年金(企業型)へ移行する際において、将来に係る勤務期間のみ
  移行する場合と過去勤務期間を含めて移行する場合のそれぞれについて、加入者の同意手続きについて簡記せよ。
\end{shortanswer}

\begin{answerbox}
  ＜将来に係る勤務期間のみ移行する場合＞

  \begin{itemize}
    \item 減額対象者の3分の1以上で組織する労働組合がある場合は、当該労働組合の同意
    \item 減額対象者の3分の2以上の個別同意（減額対象者の3分の2以上で組織する労働組合がある場合は、当該労働組合の同意で代替可能）
    \item 実施事業所が複数ある場合は、実施事業所ごとに上記の同意が必要
  \end{itemize}

  ＜過去勤務期間を含めて移行する場合＞

  
  ＜将来に係る勤務期間のみ移行する場合＞に記載の同意に加えて以下の同意も必要。
  \begin{itemize}
    \item 実施事業所ごとに移換対象の加入者の2分の1以上の同意、及び、移換対象とならない加入者の2分の1以上の同意
  \end{itemize}
\end{answerbox}


\subsubsection{給付減額の判定}

\begin{fillblank}[自作 \ 2026/10/10]

  \noindent \textbf{◯法令解釈}
  \begin{lawstructure}
    \begin{lawarticle}{}{}
      \lawparagraph{2} 
      給付の額を減額する場合の取扱い
      \lawitem{(1)}
      （略）

      \lawitem{(2)}
      次のいずれか一の場合に該当するときは、給付の額の減額として取り扱うこと。

      \hangedlawitem{\textcircled{1}}{3}
      （　　A　　）の変更によって、次のア又はイのいずれかに該当する場合

      \hangedlawitem{ア}{4}
      全部又は一部の加入者又は受給権者等に係る（　　B　　）が減少する場合

      \hangedlawitem{イ}{4}
      全部又は一部の加入者又は受給権者等に係る（　　C　　）が減少する場
      合（（　　D　　）の計算方法の変更による減少を含む。）

      \hangedlawitem{\textcircled{2}}{3}
      リスク分担型企業年金でない確定給付企業年金からリスク分担型企業年金への変更又はリスク分担型企業年金からリスク分担型企業年金でない確定給付企業年金への変更をする場合(1に該当する場合を除く。)


      \hangedlawitem{\textcircled{3}}{3}
      リスク分担型企業年金における制度変更(規則第5条第1号に規定するリスク分担型企業年金統合等変更及び規則第12条第1号に規定するリスク分担型企業年金基金合併等変更を含む。)であって、全部又は一部の加入者又は受給権者等について、積立金の額とリスク分担型企業年金掛金額の予想額の現価に相当する額を合算した額(規則第64条の規定により掛金を拠出する場合にあっては、当該拠出する額を含めるものとする。)から財政悪化リスク相当額の2分の1の額を控除した額が減少する場合(1に該当する場合を除く。) この場合において、一部の加入者又は受給権者等に係る積立金の額、リスク分担型企業年金掛金額の予想額の現価に相当する額及び財政悪化リスク相当額の算定については、通常予測給付現価により按分したものを用いること。

      \hspace*{1em}%
      ただし、加入者(受給権者を除く。)の給付設計の変更に際し、1のイに該当する場合は、（　　E　　）程度は各加入者に当該変更が行われなかったとした場合の最低積立基準額を保証する経過措置を設けており、かつ、1のアに該当しないときは、給付の額の減額として取り扱わないものとすること。また、加入者 (受給権者を除く。)の給付設計の変更に際し、1のアに該当する場合であって、その該当する加入者の（　　F　　）で組織する労働組合がある場合は、1のアに該当する各加入者の給付の名目額(基礎率のうち予定利率を（　　G　　）として算出した通常予測給付現価をいう。)が増加する給付設計の変更であり、かつ、1のイに該当しない又は（　　E　　）程度は各加入者に当該変更が行われなかったとした場合の最低積立基準額を保証する経過措置を設けるときは、あらかじめ当該労働組合の同意を得ることにより給付の額の減額として取り扱わないものとすることができること(実施事業所が2以上であるときは、全部又は一部の各実施事業所について当該同意を得ることにより、当該同意を得た実施事業所の加入者について給付の額の減額として取り扱わないものとすることができること。)。なお、通常予測給付現価又は最低積立基準額の計算に用いる基礎率は、給付設計の変更前後で同一のものを用いることとし、給付の額の算定において、規則第28 条第1項に規定する指標を用いている場合にあっては、当該指標の（　　H　　）を当該指標の見込みとして用いて計算するものとすること。

    \end{lawarticle}
  \end{lawstructure}
\end{fillblank}

\begin{answerbox}
  \begin{enumerate}[A:]
    \item 給付設計
    \item 通常予測給付現価
    \item 最低積立基準額
    \item 最低保全給付
    \item 少なくとも5年
    \item 3分の2以上
    \item 零
    \item 直近5年間の実績値の平均値
  \end{enumerate}
\end{answerbox}

\begin{explanationbox}
  令和8年9月30日発の法令解釈の改正により、給付減額とならないケース（名目額の増加）が追加されている。
\end{explanationbox}

\begin{shortanswer}[合スト \ DB編上(2026年版) \ p.163]
  解釈1-2(2)ただし書きに規定されている、名目額の増加を理由に給付減額として取り扱わない
  ことができる措置について簡記せよ。
\end{shortanswer}

\begin{answerbox}
  加入者
（受給権者を除く。）の給付設計の変更に際し、\textcircled{1}のアに該当する場合であっ
て、その該当する加入者の３分の２以上で組織する労働組合がある場合は、\textcircled{1}の
アに該当する各加入者の給付の名目額（基礎率のうち予定利率を零として算出し
た通常予測給付現価をいう。）が増加する給付設計の変更であり、かつ、\textcircled{1}のイ
に該当しない又は少なくとも５年程度は各加入者に当該変更が行われなかったと
した場合の最低積立基準額を保証する経過措置を設けるときは、あらかじめ当該
労働組合の同意を得ることにより給付の額の減額として取り扱わないものとする
ことができること（実施事業所が２以上であるときは、全部又は一部の各実施事
業所について当該同意を得ることにより、当該同意を得た実施事業所の加入者に
ついて給付の額の減額として取り扱わないものとすることができること。）。
\end{answerbox}

\begin{shortanswer}[年金1 \ H23 \ 問題2(4)\textcircled{2}（改題）]
  キャッシュバランスプランにおいて、確定給付企業年金法並びにこれに基づく政令及び省令について(法令解釈)第1-2(2)に規定する給付の額の減額判定に用いる給付現価の計算に用いる指標の予測値の算定方法について簡記せよ。
\end{shortanswer}

\begin{answerbox}
  指標の直近5年間の実績値の平均値とする。
\end{answerbox}

\newpage

\subsection{給付減額にかかる規約変更（リスク分担型DB）}

\begin{InyouJobunbox}
  \begin{itemize}
    \item 則第5条、第6条、第12条、第13条
    \item 解釈1-2(1)\textcircled{3}\textcircled{6}、解釈1-2(2)、解釈1-2(3)\textcircled{2}
  \end{itemize}
\end{InyouJobunbox}

\begin{truefalse}[年金1 \ 2018 \ 問題2(1)\textcircled{1}]

\textbf{以下の文章の正誤を判定し、誤っている場合はその理由を記述しなさい。}
\\

  リスク分担型企業年金からリスク分担型企業年金でない確定給付企業年金へ変更する場合
は、各加入者または各受給権者等の通常予測給付現価や最低積立基準額が減少しない場合で
も給付の額の減額に該当する。
\end{truefalse}

\begin{answerbox}
  \textbf{正}
\end{answerbox}

\begin{explanationbox}
  普通の給付減額になるとき：第一号〜第四号

  上記以外でリスク分担型DB実施・廃止のとき：第五号、第六号
\end{explanationbox}

\begin{truefalse}[年金1 \ H29 \ 問題2(1)]
\textbf{以下の文章の正誤を判定し、誤っている場合はその理由を記述しなさい。}
\\

  給与比例型の設計に基づくリスク分担型企業年金を実施している場合で、「一部の加入者につ
いてのみ支給率の増額の変更を行うとき」 は、給付減額にかかる同意取得手続きが必要である。
\end{truefalse}

\begin{answerbox}
  \textbf{正}
\end{answerbox}

\begin{explanationbox}
  増額の対象とならない者について最低積立基準額が低下するため、減額同意が必要
\end{explanationbox}

\begin{shortanswer}[年金1 \ 2020 \ 問題2(1)\textcircled{3}]
  リスク分担型企業年金の給付減額の理由に関する確定給付企業年金法施行規則第 5 条第5 号
の「該当することとなる蓋然性が高いこと」および同条第6 号の「該当する蓋然性が高いこ
と」として、 『確定給付企業年金法並びにこれに基づく政令及び省令について（法令解釈） 』に
記載されている内容をそれぞれ簡記しなさい。
\end{shortanswer}

\begin{answerbox}
＜規則第5条第5号＞

変更後のリスク分担型DBにおいてリスク充足額が財政悪化リスク相当額の2分の1以上であること

＜規則第5条第6号＞

変更前のリスク分担型DBにおいてリスク充足額が財政悪化リスク相当額の2分の1以下であること
\end{answerbox}


\newpage

\subsection{基金の運営}

\begin{InyouJobunbox}
  \begin{itemize}
    \item 法第18条〜第24条
    \item 令第10条の2〜第19条、第46条第2項
    \item 則第19条〜第20条
  \end{itemize}
\end{InyouJobunbox}

\begin{truefalse}[年金2 \ 2020 \ 問題1(1)\textcircled{1}]


\textbf{下線部の正誤を判定し、誤っている場合は正しい内容に改めなさい。}
\\
  
確定給付企業年金を実施する企業年金基金は、代議員会において\underline{出席した代議員の三分の
二以上}の多数により議決したとき、または基金の事業の継続が不可能となったときは、厚生
労働大臣の認可を受けて、解散することができる。
\end{truefalse}

\begin{answerbox}
  \textbf{誤}

  代議員の定数の四分の三以上
\end{answerbox}

\begin{truefalse}[年金1 \ 2023 \ 問題1(3)(ア)]
\textbf{下線部の正誤を判定し、誤っている場合は正しい内容に改めなさい。}
\\

  基金の合併に関する厚生労働大臣の認可の申請は、代議員会における代議員の定数の\underline{過半数}
による議決を経て行わなければならない。
\end{truefalse}


\begin{answerbox}
  \textbf{誤}

  四分の三以上の多数
\end{answerbox}

\begin{shortanswer}[年金2 \ 2020 \ 問題2(3)]
  「確定給付企業年金の規約の承認及び認可の基準等について」の別紙 2「確定給付企業年金の事業運営基準」において、総合型基金とは「2以上の厚生年金適用事業所の事業主が共同して実施する基金型企業年金(当該厚生年金適用事業所間の人的関係が緊密である場合等を除く。)」と定義されている。このとき、「当該厚生年金適用事業所間の人的関係が緊密である場合等」として「総合型基金のガバナンスに係る改正について」(2018(平成 30)年 9 月 27 日 付事務連絡)に挙げられている内容を簡記しなさい。
\end{shortanswer}

\begin{answerbox}
  以下のいずれかに該当する場合をいう。
  \begin{itemize}
    \item 企業の1つが行う事業と他の企業が行う事業との人間関係が緊密である場合
    \item 企業の1つが他の企業の発行済株式または出資（自己が有する自己の株式または出資を除く）の概ね２割を直接または間接に保有する関係にある場合
    \item 基金を実施する厚生年金適用事業所の全てが１つの企業に属する場合
  \end{itemize}
\end{answerbox}

\begin{explanationbox}
  過去の時事問題

  参考リンク：\url{https://www.mhlw.go.jp/stf/seisakunitsuite/bunya/0000182480.html}
\end{explanationbox}


\begin{shortanswer}[年金1 \ 2018 \ 問題2(4)]
平成30 年 4 月 1 日より、確定給付企業年金のガバナンスの改善が行われている。
総合型基金における代議員の選任のあり方について、次の(1)および(2)の内容を簡記しなさい。
\begin{enumerate}[(1)]
  \item 選定代議員の数について
  \item 選定代議員の選定に関し、選定代議員候補者の指名を希望しない事業主の選定行為の委任について
\end{enumerate}

なお、基金の設立事業主の９割以上が所属する当該基金と異なる組織体であって、通知「確定
給付企業年金の規約の承認及び認可の基準等について」に定めるものは存在しないものとする。
\end{shortanswer}

\begin{answerbox}
  \begin{enumerate}[(1)]
    \item 事業主の数の10分の1（ただし事業主の数が500を超える場合は50、30を下回る場合は3）
    \item 指名を希望しない事業主は選定行為を現役員・職員以外の第三者（選定人）に委任できる。
  \end{enumerate}
\end{answerbox}

\begin{explanationbox}
  注意：模範解答は法改正前
\end{explanationbox}

\begin{fillblank}[H29 年金1 問題1(1)]
  \noindent \textbf{◯DB法}
  \begin{lawstructure}
    \begin{lawarticle}{第十九条}{}
      次に掲げる事項は、代議員会の議決を経なければならない。
      \lawitem{一} 規約の変更
      \lawitem{二} （　　A　　）
      \lawitem{三} 毎事業年度の事業報告及び決算
      \lawitem{四} その他規約で定める事項

      \lawparagraph{2}
      代議員会は、監事に対し、基金の業務に関する監査を求め、その結果の報告を請求することができる。
    \end{lawarticle}

    \begin{lawarticle}{第七十六条}{基金の合併}
      基金は、合併しようとするときは、厚生労働大臣の認可を受けなければならない。

      \lawparagraph{2}
      前項の認可の申請は、代議員会における（　　B　　）の多数による議決を経て行わなければならない。

      \lawparagraph{3}
      （略）

      \lawparagraph{4}
      （略）
    \end{lawarticle}

    \begin{lawarticle}{第七十六条}{基金の分割}
      基金は、分割しようとするときは、厚生労働大臣の認可を受けなければならない。

      \lawparagraph{2}
      基金の分割は、（　　C　　）について行うことはできない。

      \lawparagraph{3}
      分割を行う場合においては、分割により設立される基金の加入者となるべき厚生年金保険の被保険者又は分割後存続する基金の加入者である厚生年金保険の被保険者の数が、第十二条第一項第四号（基金を共同して設立している場合にあっては、同項第五号）の政令で定める数以上であるか、又は当該数以上となることが見込まれなければならない。

      \lawparagraph{4}
      分割によって基金を設立するには、分割により設立される基金の実施事業所となるべき厚生年金適用事業所の事業主が規約を作り、その他設立に必要な行為をしなければならない。
      
      \lawparagraph{5}
      分割により設立された基金は、分割により消滅した基金又は分割後存続する基金の（　　D　　）を承継する。

      \lawparagraph{6}
      前項の規定により承継する（　　E　　）は、厚生労働大臣の認可を受けなければならない。

      \lawparagraph{7}
      前条第二項の規定は、第一項及び前項の認可の申請を行う場合について準用する。
    
    \end{lawarticle}
  \end{lawstructure}
\end{fillblank}

\begin{answerbox}
  \begin{enumerate}[A:]
    \item 毎事業年度の予算
    \item 定数の四分の三以上
    \item 実施事業所の一部
    \item 権利義務の一部
    \item 権利義務の限度
  \end{enumerate}
\end{answerbox}

\newpage

\subsection{加入者原簿}

出題例なし